\section*{\texorpdfstring{\S\CurrentExerciseGroup}{第{\CurrentExerciseGroup}节}}
\phantomsection
\addcontentsline{toc}{section}{第{\CurrentExerciseGroup}节习题}

\begin{enumerate}
\def\labelenumi{\arabic{enumi})}
\tightlist
\item
  试证：若 \(f\) 与 \(g\) 是两个数值函数 \(\geqslant 0\)（定义在 \(X\) 上），而 \(\mu\) 是 \(X\) 上的一个正测度，则
\end{enumerate}

\[
\mu^ {*} (\sup (f, g)) + \mu^ {*} (\inf (f, g)) \leqslant \mu^ {*} (f) + \mu^ {*} (g).
\]

由此推出：若 A 与 B 是 X 的任意子集，则

\[
\mu^ {*} (\mathrm{A} \cup \mathrm{B}) + \mu^ {*} (\mathrm{A} \cap \mathrm{B}) \leqslant \mu^ {*} (\mathrm{A}) + \mu^ {*} (\mathrm{B})
\]

（参见 §4, 习题 8 d)）。

\begin{enumerate}
\def\labelenumi{\arabic{enumi})}
\setcounter{enumi}{1}
\item
  对每个 \(x \in \mathbf{X}\)，试证对每个数值函数 \(f \geqslant 0\)（定义在 \(\mathbf{X}\) 上），有 \(\varepsilon_x^*(f) = f(x)\)。
\item
  在 \(\mathbf{R}\) 中给出一个不是相对紧、但其外测度（对 Lebesgue 测度而言）有限的开集的例子。
\item
  设 \(\mathrm{X}\) 是一个局部紧空间，\(\alpha\) 是一个有限数值函数 \(\geqslant 0\)（在 \(\mathrm{X}\) 上），使得对每个紧子集 \(\mathrm{K}\)（\(\mathrm{X}\) 的），\(\sum_{x\in \mathrm{K}}\alpha (x)\) 是有限的；设 \(\mu\) 是正测度
\end{enumerate}

它由质量 \(\alpha(x)\) 定义在 X 上（Ch. III, §1, No.~3, 例 I）。

\begin{enumerate}
\def\labelenumi{\alph{enumi})}
\tightlist
\item
  试证：对每个函数 \(f \geqslant 0\)（在 \(X\) 上下半连续），有 \(\mu^{*}(f) = \sum_{x \in X} \alpha(x)f(x)\)，这里约定 \(\alpha(x)f(x) = 0\) 当 \(\alpha(x) = 0\) 且 \(f(x) = -\infty\)，或者当
\end{enumerate}

\[
f (x) = + \infty
\]

\begin{enumerate}
\def\labelenumi{\alph{enumi})}
\setcounter{enumi}{1}
\tightlist
\item
  时同样成立。设 \(f \geqslant 0\) 是任一数值函数（定义在 \(X\) 上）。试证：若 \(\mu^{*}(f) < +\infty\)，则 \(\mu^{*}(f) = \sum_{x \in X} \alpha(x)f(x)\)，其中约定与 a) 中相同。（参见习题 5。）
\end{enumerate}

¶ 5) 设 X 是平面 \(\mathbf{R}^2\) 的由直线 \(D = \{0\} \times \mathbf{R}\) 与点 \((1/n, k/n^2)\) 所成之集的并组成的子集，其中 \(n\) 取遍整数 \(>0\) 的集，\(k\) 取遍有理整数集 \(\mathbf{Z}\)。

\begin{enumerate}
\def\labelenumi{\alph{enumi})}
\item
  对 D 的每个点 \((0,y)\) 及每个整数 n \textgreater{} 0，设 \(\mathrm{T}_{n}(y)\) 是 X 中点 \((u,v)\)（满足 \(u \leqslant 1/n\) 且 \(|v - y| \leqslant u\)）组成之集。试证：若对 D 中每个点 \((0,y)\) 取诸 \(\mathrm{T}_{n}(y)\) (n \textgreater{} 0) 之集作为基本邻域系，并取化为该点本身的唯一集作为 X 的其余每个点的基本邻域系，就在 X 上定义了一个拓扑 T，对它而言 X 是一个局部紧但非仿紧的空间。
\item
  设 \(\alpha\) 是 \(X\) 上的数值函数，它在 \(D\) 上等于 0，且等于 \(1/n^3\)（在每个点 \((1/n, k/n^2)\) 处）。试证：对每个紧子集 \(K\)（\(X\) 的），\(\sum_{x \in K} \alpha(x)\)
\end{enumerate}

有限；设 \(\mu\) 是 X 上由质量 \(\alpha(x)\) 定义的正测度。试证 \(\mu^{*}(D) = +\infty\)，尽管在 D 上 \(\alpha(x) = 0\)。（若关于 \(\mathcal{T}\) 的一个开集 U 包含 D，试证：在 D 上存在一个不化为一点的区间 \(a \leqslant y \leqslant b\)、一个在该区间中稠密的集 B（关于 \(\mathbf{R}\) 的通常拓扑）及一个整数 \(n > 0\)，使得对每个 \(y \in B\)，有 \(\mathrm{T}_{n}(y) \subset \mathrm{U}\)；为此可利用 Baire 定理。）

\begin{enumerate}
\def\labelenumi{\arabic{enumi})}
\setcounter{enumi}{5}
\tightlist
\item
  设 \(f\) 是一个数值函数 \(\geqslant 0\)（定义在 \(X\) 上）。
\end{enumerate}

\begin{enumerate}
\def\labelenumi{\alph{enumi})}
\item
  试证：要使 \(\mu \mapsto \mu^{*}(f)\) 从 \(\mathcal{M}_{+}(X)\) 到 \(\overline{\mathbf{R}}\) 的映射关于淡拓扑连续，必须（且只须）\(f\) 是具有紧支集的连续函数（利用习题 2）。要使 \(\mu \mapsto \mu^{*}(f)\) 关于淡拓扑下半连续，必须（且只须，参见命题 4）\(f\) 下半连续。
\item
  试证：要使 \(\mu \mapsto \mu^{*}(f)\) 从 \(\mathcal{M}_{+}(\mathrm{X})\) 到 \(\overline{\mathbf{R}}\) 的映射关于拟强拓扑（Ch. III, §1, 习题 8）连续，必须且只须 \(f\) 有界且具有紧支集（方法与 \(a\) ) 的类似）。由此推出，对每个在可数个紧集之并的余集上为零的函数 \(f \geqslant 0\)，映射 \(\mu \mapsto \mu^{*}(f)\) 关于拟强拓扑下半连续（利用定理 3）（参见习题 7 b)）。
\item
  试证：要使 \(\mu \mapsto \mu^{*}(f)\) 从 \(\mathcal{M}_{+}(\mathrm{X}) \cap \mathcal{M}^{1}(\mathrm{X})\) 到 \(\overline{\mathbf{R}}\) 的映射关于超强拓扑（Ch. III, §1, 习题 15）连续，必须且只须 \(f\) 有界；对每个函数 \(f \geqslant 0\)（定义在 \(\mathrm{X}\) 上），\(\mu \mapsto \mu^{*}(f)\) 关于超强拓扑下半连续。
\item
  试证：要使 \(\mu \mapsto \mu^{*}(f)\) 从 \(\mathcal{M}_{+}(\mathrm{X}) \cap \mathcal{M}^{1}(\mathrm{X})\) 到 \(\overline{\mathbf{R}}\) 的映射关于弱拓扑（Ch. III, §1, 习题 15）连续，必须且只须 \(f\) 在 \(\mathrm{X}\) 上连续且在无穷远点趋于 0；要使 \(\mu \mapsto \mu^{*}(f)\) 关于弱拓扑下半连续，必须且只须 \(f\) 下半连续。
\end{enumerate}

\begin{enumerate}
\def\labelenumi{\arabic{enumi})}
\setcounter{enumi}{6}
\tightlist
\item
  \begin{enumerate}
  \def\labelenumii{\alph{enumii})}
  \tightlist
  \item
    设 \((\mu_{n})\) 是局部紧空间 X 上一个递增的测度序列 \(\geqslant0\)；假设该序列在 \(\mathcal{M}_{+}(X)\) 中有上界，并用 \(\mu\) 表示它的上确界。设 f 是定义在 X 上且在可数个紧集之并的余集上为零的函数 \(\geqslant0\)。试证 \(\mu^{*}(f)=\sup\mu_{n}^{*}(f)\)（参见习题 6 b)）。
  \end{enumerate}
\end{enumerate}

\begin{enumerate}
\def\labelenumi{\alph{enumi})}
\setcounter{enumi}{1}
\tightlist
\item
  设 X 与 \(\mu\) 是习题 5 中定义的局部紧空间与测度。设 \(\alpha_{n}\) 是这样的数值函数：它由 \(\alpha\)（习题 5 的记号）在每个点 \((1/m, k/m^{2})\)（这里 \(m \leqslant n\)）处的值给出，而在 X 的其余各点处等于 0；设 \(\mu_{n}\) 是由质量 \(\alpha_{n}(x)\) 定义的测度。试证：\(\mu\) 是序列 \((\mu_{n})\) 在 \(\mathcal{M}_{+}(X)\) 中的上确界，且对一切 n 有 \(\mu_{n}^{*}(D)=0\)，但 \(\mu^{*}(D)=+\infty\)。
\end{enumerate}

\begin{enumerate}
\def\labelenumi{\arabic{enumi})}
\setcounter{enumi}{7}
\tightlist
\item
  \begin{enumerate}
  \def\labelenumii{\alph{enumii})}
  \tightlist
  \item
    设 \((\mu_n)\) 是一个递减的测度序列 \(\geqslant 0\)（在局部紧空间 \(X\) 上），并设 \(\mu\) 是该序列在 \(\mathcal{M}_+(X)\) 中的下确界。试证：若 \(f\) 是一个函数 \(\geqslant 0\)，使得从某个指标起 \(\mu_n^*(f) < +\infty\)，则 \(\mu^*(f) = \inf_n \mu_n^*(f)\)（当 \(g\) 下半连续、为正且满足 \(\mu^*(g) < +\infty\) 时，注意存在一个序列 \((h_m)\)，它由具有紧支集的连续函数 \(\geqslant 0\) 组成，使得 \(\sum_{m=1}^{\infty} h_m \leqslant g\) 且 \(\mu_n^*(g) = \sum_{m=1}^{\infty} \mu_n^*(h_m)\) 对每个指标 \(n\)。
  \end{enumerate}
\end{enumerate}

\begin{enumerate}
\def\labelenumi{\alph{enumi})}
\setcounter{enumi}{1}
\tightlist
\item
  在离散空间 X = N 上，设 \(\mu_{n}\) 是在每个点 \(m \geqslant n\) 处放置质量 +1 而定义的测度。试证：递减序列 \((\mu_{n})\) 在 \(\mathcal{M}_{+}(X)\) 中的下确界为 0，但对一切 n 有 \(\mu_{n}^{*}(X) = +\infty\)。
\end{enumerate}

\setcounter{exercisegroup}{3}
\edef\CurrentExerciseGroup{\arabic{exercisegroup}}
